\begin{textsample}{NEG Dim 1 – 2009\_c5 – Score 1.00 – 2009\_c5/enade\_2009\_publi\_15f3034279.txt}  \label{ex:f1_neg_294}
Capítulo 5 Distribuição dos Conceitos 5.1 Panorama nacional da distribuição dos conceitos A Tabela 5.1 apresenta a quantidade e o percentual de cursos de COMUNICAÇÃO SOCIAL – PUBLICIDADE E PROPAGANDA participantes do ENADE/2009 que se situaram em \textbf{cada} uma das faixas de conceitos, distribuídos de acordo com a Região em que se localizam os cursos das instituições de educação superior, mantenedoras dos cursos responsáveis pela oferta.
O Gráfico 5.1 apresenta a mesma informação sob a forma de diagrama de caixa para uma melhor visualização da dispersão dos conceitos obtidos.
No diagrama de caixa, a caixa propriamente dita fica limitada pelos quartis.
A barra dentro da caixa em negrito representa a mediana.
Linhas verticais unem os quartis aos valores extremos.
Se estes valores extremos se configurarem como valores atípicos, estes aparecem como pontos isolados e as linhas terminam no maior ( ou menor ) valor ainda não considerado atípico.
Eventuais valores atípicos aparecem como pequenos círculos com a Categoria Administrativa correspondente.
Valores considerados muito atípicos foram representados por estrelas.
Um diagrama de caixa para poucos dados ou com dados muito concentrados pode ter alguns pontos coincidentes.
Por exemplo, se mais de metade dos cursos de uma Região receberem conceito 5, o valor extremo superior, o quartil superior e a mediana coincidirão.
Se o conjunto só contém um dado, este apareceria denotado como a barra da mediana.
Se o conjunto apresenta todas as notas iguais, a notação é a mesma, somente uma barra horizontal.
Observando -se os dados da Tabela 5.1, percebe -se que, em termos nacionais, de um total de 367 cursos participantes, 45 não foram avaliados ( 12,3 \% ).
Os cursos avaliados receberam predominantemente os conceitos 3, 4 e 2, correspondendo respectivamente a 38,4 \%, 22,1 \% e a 19,6 \% do total nacional.
O conceito máximo foi atribuído a 4,1 \% dos cursos e o mínimo a 3,5 \%.
Considerando -se separadamente as cinco Regiões ( ver também Gráfico 5.1 ), observa -se que, na Região Norte, de um total de 19 cursos participantes, 6 não foram avaliados ( 31,6 \% ).
Os cursos avaliados receberam conceitos 2, mínimo e 3 correspondendo respectivamente a 52,6 \%, 10,5 \% e 5,3 \% do total nacional.
Esta Região foi a que apresentou cursos com a pior distribuição de conceitos.
Na Região Nordeste, 55 cursos participaram, dos quais, 41,8 \%, obtiveram o conceito modal 3 e 5 cursos ( 9,1 \% ) ficaram sem conceito.
O conceito máximo foi atribuído a 2 cursos ( 3,6 \% ) e o mínimo a apenas 1 curso ( 1,8 \% ).
Na Região Sudeste, 197 cursos participaram, dos quais, 43,1 \% obtiveram o conceito modal 3.
O conceito máximo foi alcançado por 6 cursos ( 3,0 \% ), 7 cursos ( 3,6 \% ) obtiveram conceito mínimo e 21 ( 10,7 \% ) cursos ficaram sem conceito.
Na Região Sul, 24 dos 62 cursos avaliados, equivalentes a 38,7 \% do total regional, classificaram -se com conceito 4.
O conceito máximo foi alcançado por 7 cursos ( 11,3 \% ), 1 curso ( 1,6 \% ) obteve conceito mínimo e 10 ( 16,1 \% ) cursos ficaram sem conceito.
Esta Região foi a que apresentou cursos com a melhor distribuição de conceitos.
Na Região Centro-Oeste, 34 cursos participaram, sendo que 44,1 \% ( 15 cursos ) obtiveram conceito 3.
O conceito máximo não foi alcançado por nenhum curso e 2 cursos ( 5,9 \% ) obtiveram conceito mínimo e 3 ( 8,8 \% ) cursos ficaram sem conceito. 5.2 Conceitos por Categoria Administrativa e por Região A Tabela 5.2 apresenta a distribuição dos cursos participantes do ENADE/2009 de COMUNICAÇÃO SOCIAL – PUBLICIDADE E PROPAGANDA, por Categoria Administrativa, de acordo com os conceitos por eles alcançados, segundo Região.
O Gráfico 5.2 apresenta a mesma informação sob a forma de diagrama de caixa para uma melhor visualização da dispersão dos conceitos obtidos.
Somente eventuais valores atípicos apareceram com o identificador da Região na qual o curso foi habilitado.
De acordo com as informações apresentadas, a grande maioria, 91,8 \%, é de cursos em instituições privadas, 4,6 \% dos cursos eram vinculados a instituições federais, 0,8 \% a instituições estaduais e os restantes 2,7 \% a instituições municipais.
As instituições estaduais foram as que apresentaram, em termos nacionais, maior proporção de cursos com conceito 5, equivalentes a 33,3 \% do total dessa categoria, mesma proporção daqueles que receberam conceito 2 e daqueles que ficaram sem conceito.
Receberam o conceito mínimo, em instituições federais, 3 cursos ( 17,6 \% ), mesmo número e proporção daqueles que obtiveram conceitos 3 e 4.
Ficaram sem conceito 4 cursos ( 23,5 \% ) em instituições federais.
Entre as instituições municipais, o conceito modal foi 3 e 4 ( 40,0 \% \textbf{cada} ) seguido dos conceitos 1 e 2, com 10,0 \% \textbf{cada}.
Nenhum curso recebeu o conceito máximo nem ficou sem conceito nas instituições municipais.
Os cursos nas instituições privadas tiveram como conceito modal o valor 3, com uma incidência de 39,8 \%.
Somente 10 cursos ( 3,0 \% ) receberam o conceito máximo, 9 ( 2,7 \% ) receberam o conceito mínimo e 40 ( 11,9 \% ) ficaram sem conceito. com 10,0 \% \textbf{cada}.
Nenhum curso recebeu o conceito máximo nem ficou sem conceito nas instituições municipais.
Os cursos nas instituições privadas tiveram como conceito modal o valor 3, com uma incidência de 39,8 \%.
Somente 10 cursos ( 3,0 \% ) receberam o conceito máximo, 9 ( 2,7 \% ) receberam o conceito mínimo e 40 ( 11,9 \% ) ficaram sem conceito.
Em todas as Regiões as instituições privadas constituíram a grande maioria dos cursos de COMUNICAÇÃO SOCIAL – PUBLICIDADE E PROPAGANDA participantes.
Na Região Norte, dos 17 cursos participantes em instituições privadas 5 ficaram sem conceito, 10 obtiveram conceito 2, 1 o conceito 3 e 1 o conceito mínimo.
Não houve participação de cursos em instituições estaduais nem em instituição municipal.
Nas instituições federais, dos 2 cursos participantes, 1 ficou sem conceito e o outro recebeu o conceito mínimo.
Na Região Nordeste, dos 4 cursos participantes em instituições federais 2 ficaram sem conceito, 1 com conceito máximo e o outro com o conceito 4.
O único curso participante em instituição estadual ficou com conceito 2.
Não houve participação de curso de instituição municipal.
Nas instituições privadas, dos 50 cursos participantes, 3 ficaram sem conceito, 23 com conceito 3, 13 com conceito 4 e 9 com conceito 2.
Os conceitos mínimo e máximo foram obtidos por 1 curso \textbf{cada}.
Na Região Sudeste, dos 4 cursos participantes em instituições federais, 1 obteve o conceito mínimo, 1 o conceito máximo, 1 o conceito 3 e outro o conceito 4.
O único curso participante em instituição estadual ficou sem conceito.
Dos 7 cursos de instituição municipal, 3 obtiveram conceito 3, 2 o conceito 4, 1 o conceito 2 e 1 o conceito mínimo.
Nenhum curso em instituição municipal ficou sem conceito.
Nas instituições privadas, dos 185 cursos participantes, 20 ficaram sem conceito, 81 com conceito 3, 36 com conceito 4 e 38 com conceito 2.
Os conceitos mínimo e máximo foram obtidos por 5 cursos \textbf{cada}.
Na Região Sul, dos 4 cursos participantes em instituições federais, apenas 1 ficou sem conceito, 1 obteve o conceito mínimo e os outros 2 o conceito máximo.
O único curso participante em instituição estadual, ficou com conceito 5.
Dos 3 cursos de instituições municipais, nenhum ficou sem conceito, 2 obtiveram conceito 4 e 1, o conceito 3.
Nas instituições privadas, dos 54 cursos participantes, 9 ficaram sem conceito, 22 com conceito 4, 16 com conceito 3 e 3 com conceito 2.
O conceito máximo foi obtido por 4 cursos e nenhum curso obteve o conceito mínimo.
Na Região Centro-Oeste, dos 3 cursos participantes em instituições federais 2 ficaram com conceito 3 e 1 com conceito 4.
Não houve participação de curso em instituição estadual nem em instituição municipal.
Nas instituições privadas, dos 31 cursos participantes, 3 ficaram sem conceito, 13 com conceito 3, 10 com conceito 2 e 3 com conceito 4.
O conceito mínimo foi obtido por 2 cursos e nenhum obteve o conceito máximo. 5.3 Conceitos por Organização Acadêmica e por Região Na Tabela 5.3 encontra -se a distribuição dos conceitos atribuídos aos cursos participantes do ENADE/2009 de COMUNICAÇÃO SOCIAL – PUBLICIDADE E PROPAGANDA, por Organização Acadêmica, segundo as Regiões.
O Gráfico 5.3 apresenta a mesma informação sob a forma de diagrama de caixa para uma melhor visualização da dispersão dos conceitos obtidos.
Somente eventuais valores atípicos apareceram com o identificador da Categoria Administrativa.
Os cursos desta área \textbf{estão} principalmente alocados em Faculdades ( 47,7 \% ), depois em Universidades ( 33,2 \% ) e Centros Universitários ( 18,8 \% ).
Este ordenamento se repete em todas as Regiões, com exceção da Região Norte, na qual Centros Universitários \textbf{estão} em segundo lugar.
A presença em Institutos Federais de Educação, Ciência e Tecnologia ( IFECT ) é incipiente, com apenas 1 curso participante, sendo este na Região Norte, que ficou sem conceito.
Não houve participação de Centros Federais de Educação Tecnológica ( CEFET ).
Nas Universidades, Centros Universitários e Faculdades, o conceito modal foi 3 ( 36,9 \%, 42,0 \% e 38,3 \%, respectivamente ).
Proporcionalmente poucos cursos nos 3 tipos de Organização Acadêmica mais frequentes, Universidades, Centros Universitários e Faculdades, receberam conceito mínimo, respectivamente 4,1 \%, 4,3 \% e 2,9 \%.
Uma proporção também reduzida de cursos atingiu o conceito máximo : 7,4 \% em Universidades, 4,3 \% em Centros Universitários e 1,7 \% em Faculdades.
Nos 3 tipos de Organização Acadêmica existiram cursos que ficaram sem conceito ( respectivamente 4,9 \%, 13,0 \% e 16,6 \% ).
Considerando -se separadamente as cinco Regiões, observou -se que, na Região Norte não apresentou nenhum curso em CEFET e somente um curso em IFECT, que ficou sem conceito.
Dos 2 cursos em Universidades, 1 com conceito 2 e outro com conceito 1.
Dos 4 cursos em Centros Universitários, 1 ficou sem conceito, 1 com conceito 2, 1 com conceito 3 e 1 com conceito mínimo.
Nas Faculdades, dos 12 cursos participantes, 4 ficaram sem conceito e os demais com conceito 2.
Como já comentado, nenhum curso nesta Região obteve conceito máximo nem conceito mínimo.
Na Região Nordeste, dos 12 cursos em Universidades, 2 ficaram sem conceito, 4 com conceito 4, 2 com conceito 3, 3 com conceito 2 e 1 com conceito máximo.
Dos 3 cursos em Centros Universitários, 2 ficaram com conceito 2 e 1 com conceito 4.
Nas Faculdades, dos 40 cursos participantes, 3 ficaram sem conceito, 21 com conceito 3, 9 com conceito 4 e 5 com conceito 2.
Os conceitos máximo e mínimo foram obtidos por 1 curso \textbf{cada}.
Na Região Sudeste, dos 73 cursos em Universidades, 3 ficaram sem conceito, 31 com conceito 3, 18 com conceito 4 e 15 com conceito 2.
Nas universidades, os conceitos máximo e mínimo foram obtidos por 3 cursos \textbf{cada}.
Dos 46 cursos em Centros Universitários, 6 ficaram sem conceito, 21 com conceito 3, 10 com conceito 2, 7 com conceito 4, 1 com conceito mínimo e 1 com conceito máximo.
Nas Faculdades, dos 78 cursos participantes, 12 ficaram sem conceito, 33 com conceito 3, 14 com conceito 2 e 14 com conceito 4.
Nestas instituições, o conceito máximo foi obtido por 2 cursos e 3 ficaram com conceito mínimo.
Na Região Sul dos 25 cursos em Universidades, 1 ficou sem conceito, 11 com conceito 4, 6 com conceito 3 e 1 com conceito 2.
Nas universidades, os conceitos máximo e mínimo foram obtidos por 5 e 1 cursos, respectivamente.
Dos 10 cursos em Centros Universitários, apenas 1 ficou sem conceito, 4 com conceito 3 e 3 com conceito 4.
Nos Centros Universitários, o conceito máximo foi obtido por 2 cursos.
Nas Faculdades, dos 27 cursos participantes, 8 ficaram sem conceito, 10 com conceito 4, 7 com conceito 3 e 2 com conceito 2.
Nestas instituições, o conceito máximo e o conceito mínimo não foram atribuídos a nenhum curso.
Na Região Centro-Oeste dos 10 cursos em Universidades, nenhum curso ficou sem conceito, 6 com conceito 3, 2 com conceito 2 e 2 com conceito 4.
Nas universidades, os conceitos máximo e mínimo não foram obtidos por nenhum curso.
Dos 6 cursos em Centros Universitários, 1 ficou sem conceito, 3 com conceito 3, 1 com conceito 2.
Nos Centros Universitários, o conceito mínimo foi obtido por 1 curso e o máximo por nenhum.
Nas Faculdades, dos 18 cursos participantes, 2 ficaram sem conceito, 7 com conceito 2, 6 com conceito 3 e 2 com conceito 4.
Nestas instituições, o conceito mínimo foi atribuído a 1 único curso e nenhum curso alcançou o conceito máximo.

% matched lemmas: cada, estar
\end{textsample}
